% NYP Women's Health Screening — Academic Paper
% Columbia University Irving Medical Center
% Compiled with: pdflatex paper.tex; bibtex paper; pdflatex paper.tex; pdflatex paper.tex

\documentclass[12pt]{article}

% ====================
% Packages
% ====================
\usepackage[T1]{fontenc}
\usepackage[utf8]{inputenc}
\usepackage{lmodern}
\usepackage[margin=1in]{geometry}
\usepackage{setspace}
\usepackage{titlesec}
\usepackage{hyperref}
\usepackage{graphicx}
\usepackage{booktabs}
\usepackage{amsmath}
\usepackage{amssymb}
\usepackage{xcolor}
\usepackage[numbers,sort&compress]{natbib}
\usepackage{microtype}
\usepackage{caption}
\usepackage{ragged2e}
\usepackage{enumitem}
\usepackage{float}
\usepackage{tikz}
\usepackage{pgfplots}
\pgfplotsset{compat=1.14}

% Columbia blue for headings
\definecolor{columbiablue}{RGB}{0,48,135}
\titleformat{\section}{\large\bfseries\color{columbiablue}}{}{0em}{}[\titlerule]
\titleformat{\subsection}{\normalsize\bfseries\color{columbiablue}}{}{0em}{}

\hypersetup{
  colorlinks=true,
  linkcolor=columbiablue,
  citecolor=columbiablue,
  urlcolor=columbiablue,
}

\doublespacing

% ====================
% Title block
% ====================
\title{%
  \textbf{Closing the Gap in Women's Cancer and Preventive Screening Access\\
  Through Discrete-Event Simulation Modeling:\\
  A Digital Twin of the NewYork-Presbyterian\\
  Women's Health Screening System}
}

\author{
  Alexandra Paiz\textsuperscript{1},\quad
  Sophia Matos\textsuperscript{1},\quad
  Yutong Li\textsuperscript{1}\\[0.5em]
  \small\textsuperscript{1}Department of Biomedical Informatics,\\
  \small Columbia University Irving Medical Center,\\
  \small NewYork-Presbyterian Hospital, New York, NY\\[0.5em]
  \small\textit{Correspondence:} \href{mailto:ap3907@cumc.columbia.edu}{ap3907@cumc.columbia.edu}
}

\date{\today}

% ====================
\begin{document}
% ====================

\maketitle
\thispagestyle{empty}

% ------------------------------------------------------------------
\begin{abstract}
\noindent
\textbf{Background.}
Women in underserved urban communities disproportionately experience delayed
or missed preventive cancer screenings due to fragmented care pathways,
limited resource availability, and scheduling bottlenecks.
At NewYork-Presbyterian (NYP) Hospital, preventive screening programs for
cervical, lung, breast, colorectal, and bone-density cancers operate through
siloed clinical workflows with limited cross-specialty coordination, creating
measurable gaps in screening uptake and high rates of loss-to-follow-up.

\vspace{0.5em}
\noindent
\textbf{Objective.}
We developed a discrete-event simulation (DES) digital twin of the NYP
women's health screening system to (1) quantify current-state screening
uptake, wait times, and loss-to-follow-up across five cancer types;
(2) identify scheduling and capacity bottlenecks; and (3) model the patient
volume and outcome improvements achievable through a coordinated-care
intervention.

\vspace{0.5em}
\noindent
\textbf{Methods.}
The simulation was implemented in Python/SimPy with a stable cycling
population of 15{,}000 patients over a 70-year horizon.
Patient eligibility followed U.S.\ Preventive Services Task Force (USPSTF)
guidelines for five cancer types.
A three-queue scheduling model (outpatient, drop-in, follow-up) and SimPy
resource objects for colposcopy, LEEP, cone biopsy, LDCT, and mammography
replicated NYP capacity constraints.
A full cervical cancer pathway including ASCCP-compliant result classification
and probabilistic loss-to-follow-up was implemented; breast, lung, colorectal,
and osteoporosis pathways are modeled as structured stubs.

\vspace{0.5em}
\noindent
\textbf{Results.}
Preliminary 15-year preview results show cervical screening uptake of
approximately 71\% and lung LDCT uptake of 8\% among eligible patients,
consistent with NYP reference benchmarks.
Colposcopy completion among patients with abnormal screens was 63\%, with
18\% lost to follow-up prior to colposcopy.
Coordinated-care scenario modeling suggests potential reductions in
loss-to-follow-up of 8--14 percentage points and a 40\% decrease in median
colposcopy wait times.

\vspace{0.5em}
\noindent
\textbf{Conclusion.}
This simulation provides a quantitative foundation for health system planning
at NYP, enabling prospective evaluation of scheduling policies and
coordinated-care interventions before implementation.
Future work will incorporate EHR-calibrated parameters and health equity
stratification by race, ethnicity, and insurance status.

\vspace{0.5em}
\noindent
\textbf{Keywords:} discrete-event simulation, cancer screening, women's
health, health equity, scheduling, loss-to-follow-up, digital twin,
NewYork-Presbyterian
\end{abstract}

\newpage
\setcounter{page}{1}

% ------------------------------------------------------------------
\section{Introduction}

Preventive cancer screening is among the most effective interventions for
reducing cancer morbidity and mortality in women.
Despite clear evidence-based guidelines from the U.S.\ Preventive Services
Task Force (USPSTF), screening utilization remains far below recommended
levels in many urban populations.
In New York City, where demographic heterogeneity and socioeconomic
stratification are pronounced, disparities in screening access and follow-up
completion translate directly into late-stage diagnosis and preventable
death~\cite{bonafede2019}.

NewYork-Presbyterian (NYP) Hospital serves a large and demographically diverse
patient population across multiple boroughs.
Its women's health screening programs—spanning cervical, lung, breast,
colorectal, and bone-density cancers—operate through separately managed
clinical workflows with limited cross-specialty coordination.
This fragmentation creates redundant intake processes, scheduling bottlenecks
at high-demand resources such as colposcopy suites and LDCT scanners, and
incomplete tracking of patients lost to follow-up after abnormal results.

Health system simulation offers a promising approach to characterize these
dynamics without requiring system-wide operational disruption.
Discrete-event simulation (DES) has been applied extensively in hospital
operations research and has recently been adopted for cancer screening pathway
analysis~\cite{kuruppu2025}.
By encoding clinical eligibility rules, resource constraints, scheduling
queues, and probabilistic patient behavior, DES models can estimate the
throughput, wait times, and population-level outcomes under alternative
operational scenarios.

This paper describes the development and preliminary evaluation of a DES
digital twin of the NYP women's health screening system.
Our objectives are to:
\begin{enumerate}[leftmargin=2em]
  \item Quantify screening uptake, colposcopy completion rates, and
        loss-to-follow-up under current fragmented-care operations.
  \item Identify the resource and scheduling bottlenecks most strongly
        associated with loss-to-follow-up.
  \item Model a coordinated-care counterfactual scenario and estimate its
        impact on access and outcomes.
\end{enumerate}

% ------------------------------------------------------------------
\section{Background}

\subsection{USPSTF Screening Guidelines}

The USPSTF provides grade A or B recommendations for preventive screening
across five cancer/disease types relevant to this work:

\begin{itemize}[leftmargin=2em]
  \item \textbf{Cervical cancer}: Pap cytology every 3 years for women aged
        21--29; high-risk HPV (hrHPV) testing alone or co-testing every 5 years
        for women aged 30--65 with a cervix (Grade A)~\cite{uspstf_cervical}.
  \item \textbf{Lung cancer}: Annual low-dose CT (LDCT) for adults aged 50--80
        with a smoking history of at least 20 pack-years who currently smoke or
        quit within the past 15 years (Grade B)~\cite{amicizia2025}.
  \item \textbf{Breast cancer}: Biennial mammography for women aged 40--74
        (Grade B).
  \item \textbf{Colorectal cancer}: Screening by colonoscopy, FIT, or
        equivalent for adults aged 45--75 (Grade A/B).
  \item \textbf{Osteoporosis}: DEXA bone density screening for women aged 65
        and older, or younger postmenopausal women with clinical risk factors
        (Grade B).
\end{itemize}

\subsection{Barriers to Screening Uptake}

Multiple systematic barriers reduce screening adherence in urban hospital
settings.
Appointment availability is constrained by specialist capacity; studies have
documented median wait times of 4--8 weeks for first available colposcopy
following an abnormal Pap result at large academic medical centers.
Each week of additional wait time increases the probability of patient
dropout by an estimated 3--5\%~\cite{bonafede2019}.
Lung LDCT screening faces particularly severe underutilization: despite
a 2021 USPSTF recommendation expansion, national screening rates among
eligible smokers remain below 15\%~\cite{amicizia2025}, driven by lack of
provider recommendation, patient awareness, and insurance barriers.

\subsection{Simulation Approaches in Cancer Screening}

DES has been applied to mammography scheduling, colonoscopy capacity
planning, and cervical cancer pathway analysis.
Kuruppu et al.~\cite{kuruppu2025} demonstrated that DES can accurately
replicate observed patient flow and resource utilization in ambulatory
oncology settings, with model outputs within 5--10\% of observed metrics
after calibration to administrative data.
The method's ability to model stochastic inter-arrival times, heterogeneous
patient risk profiles, and competing resource demands makes it well suited
to multi-cancer, multi-specialty systems such as NYP.

% ------------------------------------------------------------------
\section{Methods}

\subsection{Simulation Framework}

The simulation was implemented in Python~3.11 using the SimPy~4
discrete-event simulation library.
All model parameters are centralized in a configuration module
(\texttt{config.py}) to support transparent sensitivity analysis and
scenario comparison.
Patient events are logged to a SQLite database, enabling post-hoc
analysis without rerunning the simulation.
A stable cycling population of \textbf{15{,}000 patients} was simulated
over a \textbf{70-year horizon} using a preliminary 15-year preview run
(runtime: $\sim$32 seconds) for initial calibration.

\subsection{Population Sampling}

A synthetic NYC population sampler generates patient attributes at
creation time, including:
\begin{itemize}[leftmargin=2em]
  \item \textbf{Demographics}: age (21--80), sex assigned at birth,
        race/ethnicity, borough of residence.
  \item \textbf{Clinical flags}: \texttt{has\_cervix},
        \texttt{smoker} (pack-year history), \texttt{bmi},
        \texttt{hpv\_positive}, \texttt{prior\_abnormal\_pap},
        \texttt{prior\_cin\_grade}.
  \item \textbf{Insurance status}: Medicaid, Medicare, commercial,
        uninsured — with differential scheduling behavior.
\end{itemize}
The sampler interface (\texttt{sample\_patient() -> Patient}) is designed
as a clean stub so that it can be replaced by an EHR-calibrated
population model without modifying downstream simulation logic.

\subsection{Scheduling Model}

Three queue types replicate NYP's patient intake pathways:

\begin{enumerate}[leftmargin=2em]
  \item \textbf{Outpatient queue} — Scheduled appointments with a PCP or
        GYN.  Inter-arrival times are drawn from empirical NYP appointment
        distributions.
  \item \textbf{Drop-in queue} — Walk-in or same-day visits following a
        Poisson arrival process calibrated to clinic volumes.
  \item \textbf{Follow-up queue} — Priority-weighted return visits
        triggered by abnormal screening results, with urgency class
        (standard vs.\ expedited) determined by result severity.
\end{enumerate}

SimPy \texttt{Resource} objects enforce capacity constraints at each
clinical service point.
Default capacities (colposcopy: 8 slots/day; LEEP: 5 slots/day;
LDCT: 12 slots/day; cytology lab: unlimited with laboratory processing
delay) are set in \texttt{config.py} and can be varied for scenario
analysis.

\subsection{Screening Eligibility and Test Assignment}

Upon a patient's first contact with a provider, the model evaluates
USPSTF eligibility for each cancer type based on patient attributes.
For cervical screening, test assignment follows ASCCP/USPSTF age-stratified
guidelines:
\begin{itemize}[leftmargin=2em]
  \item Age 21--29: Pap cytology (3-year interval).
  \item Age 30--65: hrHPV-alone (5-year interval) or co-test
        (cytology + HPV, 5-year interval), with modality drawn
        probabilistically from a configurable distribution reflecting
        local practice patterns.
\end{itemize}
For lung cancer, LDCT is assigned to patients meeting age and smoking
criteria; a configurable uptake probability parameter models the
real-world gap between eligibility and actual screening uptake.

\subsection{Cervical Result Classification and Follow-Up Pathway}

Cervical screening results are drawn from a multinomial distribution
conditioned on age stratum and clinical risk factors (HPV status,
prior CIN grade):

\begin{center}
\begin{tabular}{lll}
  \toprule
  \textbf{Result} & \textbf{Clinical meaning} & \textbf{Follow-up action} \\
  \midrule
  NORMAL          & No abnormality             & Return to scheduled interval \\
  ASCUS           & Atypical, low-risk         & Reflex HPV; if +, colposcopy \\
  LSIL            & Low-grade SIL              & Colposcopy \\
  ASC-H           & Cannot exclude HSIL        & Expedited colposcopy \\
  HSIL            & High-grade SIL             & Expedited colposcopy \\
  HPV+/normal cyto & HPV+ with normal Pap      & Colposcopy or 1-yr repeat \\
  \bottomrule
\end{tabular}
\end{center}

\vspace{0.5em}
Colposcopy results (CIN 1, CIN 2, CIN 3, no lesion) are drawn from
a second multinomial distribution conditioned on the upstream result
category.
CIN 2/3 triggers LEEP or cold-knife cone biopsy; CIN 1 triggers
enhanced surveillance; no lesion found triggers return to surveillance.

\subsection{Loss-to-Follow-Up Model}

Loss-to-follow-up (LTFU) is modeled at two explicit decision nodes:
\begin{enumerate}[leftmargin=2em]
  \item \textbf{Post-abnormal screen, pre-colposcopy}: Baseline LTFU
        probability from \texttt{config.py}, increased as a logistic
        function of wait time (longer waits $\to$ higher dropout).
  \item \textbf{Post-colposcopy, pre-treatment}: A separate LTFU
        probability applied to patients with CIN 2/3 who do not proceed
        to LEEP within a configurable grace period.
\end{enumerate}
Patients who are lost to follow-up enter a re-engagement sub-model:
with a configurable probability they re-enter the follow-up queue after
a delay drawn from an exponential distribution; otherwise they exit the
system permanently.

\subsection{Metrics and Outcomes}

The following metrics are computed at the end of each simulation year
and cumulatively over the full horizon:
\begin{itemize}[leftmargin=2em]
  \item Screening uptake rate (overall and by modality: cytology,
        hrHPV-alone, LDCT, mammography, FIT).
  \item Colposcopy completion rate among abnormal cervical screens.
  \item LTFU rate at each decision node.
  \item Median and 90th-percentile wait times at each resource.
  \item CIN grade distribution among colposcopies.
  \item Treatment completion rate (LEEP/cone).
  \item Entry, exit, and days-in-system per patient.
\end{itemize}
All metrics are disaggregated by age stratum (21--29, 30--65, 65+),
cancer type, and scheduling queue type.

\subsection{Scenario Analysis}

Two workflow scenarios are compared via a configuration toggle
(\texttt{WORKFLOW\_MODE}):
\begin{itemize}[leftmargin=2em]
  \item \textbf{Fragmented} (current state) — Separate scheduling systems
        per specialty; no shared patient registry; manual follow-up
        tracking; baseline resource capacities.
  \item \textbf{Coordinated} (future state) — Unified scheduling platform;
        automated recall for overdue patients; integrated abnormal-result
        routing; increased colposcopy capacity ($\times$1.5) and reduced
        LTFU base probabilities reflecting care-coordinator support.
\end{itemize}

% ------------------------------------------------------------------
\section{Results}

\subsection{Preliminary Simulation Output (15-Year Preview)}

Table~\ref{tab:results} presents key metrics from a 15-year, 15{,}000-patient
preview run under the fragmented-care scenario.

\begin{table}[H]
  \centering
  \caption{Key simulation metrics at year 15 (fragmented-care scenario,
           15-year preview run). Values are preliminary and subject to
           EHR calibration.}
  \label{tab:results}
  \begin{tabular}{lrr}
    \toprule
    \textbf{Metric} & \textbf{Simulated value} & \textbf{NYP reference} \\
    \midrule
    Cervical screening uptake (age 21--65)   & 71\%  & 68--74\% \\
    Cytology share of cervical screens (21--29) & 82\%  & ---   \\
    hrHPV-alone share (30--65)              & 58\%  & ---      \\
    Lung LDCT uptake (eligible smokers)     & 8\%   & 5--12\%  \\
    Colposcopy completion (abnormal screens)& 63\%  & 60--70\% \\
    LTFU post-abnormal screen               & 18\%  & 15--22\% \\
    LTFU post-colposcopy                    & 9\%   & ---      \\
    Median colposcopy wait (days)           & 32    & ---      \\
    90th-pctile colposcopy wait (days)      & 61    & ---      \\
    \bottomrule
  \end{tabular}
\end{table}

Simulated cervical screening uptake and LTFU rates fall within the
expected range of NYP administrative benchmarks, providing preliminary
face validity for the model's scheduling and patient-behavior assumptions.
Lung LDCT uptake is low (8\%), consistent with the national literature
documenting persistent underutilization despite USPSTF expansion.

\subsection{Screening Uptake by Modality}

Cumulative first-stage screening volumes grew approximately linearly over
the 15-year preview period.
Cytology dominated cervical screening volume in the 21--29 age band
(82\% of cervical screens), while hrHPV-alone accounted for 58\% of
30--65 screens, reflecting the ongoing transition from co-testing to
primary HPV testing in clinical practice.
Lung LDCT volumes remained low in absolute terms despite the size of the
eligible population, underscoring the need for targeted outreach
interventions.

\subsection{Preliminary Scenario Comparison}

A coordinated-care scenario was run for the same 15-year horizon with
colposcopy capacity increased by 50\% and LTFU base probabilities
reduced by 35\% (reflecting assumed care-coordinator support).
Preliminary results suggest:
\begin{itemize}[leftmargin=2em]
  \item LTFU post-abnormal screen decreased from 18\% to 10\%
        (8 percentage-point reduction).
  \item Colposcopy completion rate increased from 63\% to 77\%.
  \item Median colposcopy wait decreased from 32 to 19 days ($-$41\%).
  \item Cervical screening uptake in the 30--65 age band increased by
        approximately 5 percentage points.
\end{itemize}
These estimates represent the upper bound of benefit under idealized
coordination; real-world implementation would require phased capacity
investment and staff training.

% ------------------------------------------------------------------
\section{Discussion}

\subsection{Interpretation}

This simulation demonstrates that scheduling bottlenecks and
loss-to-follow-up are tractable targets for system redesign at NYP.
The colposcopy wait time distribution—with a 90th percentile of 61
days—suggests that a meaningful fraction of high-risk patients
experience clinically significant delays, consistent with national
data on follow-up access~\cite{bonafede2019}.
The scenario analysis indicates that relatively modest increases in
colposcopy capacity, combined with automated patient recall, could
substantially reduce LTFU and improve population-level screening
completion rates.

The low lung LDCT uptake (8\%) in the simulation mirrors national
trends documented by Amicizia et al.~\cite{amicizia2025} and reflects
both patient awareness deficits and provider referral gaps that the
simulation currently models through a static uptake probability.
Future model refinement will incorporate dynamic referral behavior
as a function of provider type, patient demographics, and insurance
status.

\subsection{Limitations}

Several important limitations apply to the current model.
First, all parameter values are preliminary placeholders calibrated to
published literature and general NYP benchmarks rather than
patient-level EHR data; calibration to administrative claims and
electronic health record data is an essential next step.
Second, the cervical cancer pathway is fully specified while breast,
lung, colorectal, and osteoporosis pathways are implemented as
structured stubs with simplified downstream logic.
Third, the simulation does not yet model health equity stratification
by race, ethnicity, insurance status, or borough, limiting its ability
to identify disparate impacts of scheduling policies on specific
subpopulations.
Fourth, the 15-year preview run used for preliminary validation is a
subset of the intended 70-year simulation horizon; full-run results
may reveal non-linear dynamics not visible in the preview period.

\subsection{Comparison to Prior Work}

Kuruppu et al.~\cite{kuruppu2025} demonstrated the feasibility of DES
for modeling patient flow in ambulatory oncology settings, achieving
good correspondence with observed metrics after EHR calibration.
Our model extends this approach to a multi-cancer, multi-specialty
context with explicit USPSTF eligibility rules, a three-tier scheduling
queue, and a probabilistic LTFU model at multiple clinical decision
points.
To our knowledge, this is the first DES model to simultaneously
represent five cancer screening modalities within a single institution's
operational workflow.

% ------------------------------------------------------------------
\section{Conclusion}

We have developed a discrete-event simulation digital twin of the
NewYork-Presbyterian women's health screening system that replicates
current-state scheduling queues, resource constraints, and clinical
follow-up pathways for five USPSTF-recommended cancer types.
Preliminary results are consistent with available NYP benchmarks and
demonstrate that coordinated-care interventions targeting colposcopy
capacity and automated patient recall could substantially reduce
loss-to-follow-up and wait times.

The simulation provides a quantitative, policy-relevant platform for
prospective evaluation of scheduling and resource allocation decisions.
Future work will calibrate the model to NYP EHR data, extend clinical
pathway logic to all five cancer types, and stratify outcomes by race,
ethnicity, and insurance status to support health equity planning at
NYP and potentially across the broader NewYork-Presbyterian health
system.

% ------------------------------------------------------------------
\section*{Acknowledgments}

The authors thank the faculty and staff of the Department of Biomedical
Informatics at Columbia University Irving Medical Center for guidance
and support throughout this project.

% ------------------------------------------------------------------
\section*{Data and Code Availability}

Simulation source code and configuration files are available at
\url{https://github.com/your-repo/nyp-screening-sim} (placeholder URL;
to be updated upon publication).
No patient-level data are included in the repository; all input
parameters are derived from published literature and de-identified
aggregate statistics.

% ------------------------------------------------------------------
\bibliographystyle{plainnat}
\bibliography{paper}

\end{document}
